% ═════════════════════════════════════════════════════════════
% A1 — Limited-Memory BFGS (L-BFGS)
\label{sec:lbfgs}
% ═════════════════════════════════════════════════════════════


% ─────────────────────────────────────────────────────────────
\subsection{Quasi-Newton Methods: Background}
% ─────────────────────────────────────────────────────────────

\noindent Quasi-Newton methods are iterative algorithms for unconstrained
optimization that approximate the Newton direction without explicitly computing
the true Hessian matrix~\cite{nocedal2006numerical}. Recall that the pure Newton
step would be $p_k = -[\nabla^2 f(w_k)]^{-1} \nabla f(w_k)$: it accounts for
the local curvature and converges quadratically near the solution, but requires
computing and factorizing an $m \times m$ matrix at every iteration, which costs
$O(m^3)$ and is infeasible for large $m$. Quasi-Newton methods replace the true
inverse Hessian with an approximation $H_k \approx [\nabla^2 f(w_k)]^{-1}$ that
is built cheaply from first-order information only. The search direction is then:
\begin{equation}
\label{eq:qn_direction}
    p_k = -H_k \nabla f(w_k),
\end{equation}
and the iterate is updated as $w_{k+1} = w_k + \alpha_k p_k$, where $\alpha_k > 0$
is the step length determined by a line search.

\paragraph{The secant equation.}
After each step, the algorithm observes two key vectors: the parameter displacement
$s_k = w_{k+1} - w_k$ and the gradient difference
$y_k = \nabla f(w_{k+1}) - \nabla f(w_k)$. These encode curvature information:
in one dimension, the second derivative can be approximated as
$f''(x) \approx \Delta f' / \Delta x = y_k / s_k$. In $\mathbb{R}^m$, this
generalizes to the requirement that the new approximation $H_{k+1}$ maps the
gradient change $y_k$ to the step $s_k$, exactly as the true inverse Hessian
would near a quadratic:
\begin{equation}
\label{eq:secant}
    H_{k+1}\, y_k = s_k.
\end{equation}
This is called the \emph{secant condition}. Geometrically, it forces $H_{k+1}$
to reproduce the correct curvature along the direction just explored. However,
in dimension $m > 1$, equation~\eqref{eq:secant} has infinitely many solutions:
it constrains $H_{k+1}$ only along $y_k$, leaving all orthogonal directions
undetermined.

\paragraph{The BFGS update.}
The BFGS method resolves the non-uniqueness by solving a
constrained optimization problem: among all symmetric positive definite matrices
satisfying the secant condition, find the one closest to the previous
approximation $H_k$ in a weighted Frobenius norm. This \emph{minimal-change}
principle ensures that information accumulated in $H_k$ is preserved as much as
possible, and only updated in the directions where new curvature has been
observed. The unique solution is:
\begin{equation}
\label{eq:bfgs}
    H_{k+1} = V_k^T H_k V_k + \rho_k s_k s_k^T,
\end{equation}
where
\begin{equation}
    \rho_k = \frac{1}{y_k^T s_k}, \qquad V_k = I - \rho_k y_k s_k^T.
\end{equation}

\noindent The structure of this update is highly deliberate. The scalar $\rho_k$
normalizes the curvature so that $\rho_k(s_k^T y_k) = 1$. The matrix $V_k$ acts
as a projector that annihilates $y_k$: since $V_k y_k = y_k - \rho_k y_k
(s_k^T y_k) = y_k - y_k = 0$, the first term in~\eqref{eq:bfgs} contributes
nothing along $y_k$. The second term $\rho_k s_k s_k^T$ then injects precisely
the curvature needed to satisfy the secant condition. Indeed, multiplying
\eqref{eq:bfgs} by $y_k$:
\begin{equation}
    H_{k+1} y_k
    = V_k^T H_k \underbrace{(V_k y_k)}_{=\,0}
      + \rho_k s_k \underbrace{(s_k^T y_k)}_{=\,1/\rho_k}
    = s_k. \checkmark
\end{equation}

\noindent While BFGS achieves superlinear convergence on strongly convex problems,
it has a fundamental limitation: after $k$ iterations, $H_k$ is a dense
$m \times m$ matrix. Storing it requires $O(m^2)$ memory, and updating it via
formula~\eqref{eq:bfgs} costs $O(m^2)$ per iteration. For large $m$, both
become prohibitive.


% ─────────────────────────────────────────────────────────────
\subsection{The L-BFGS Algorithm}
% ─────────────────────────────────────────────────────────────

\noindent The Limited-Memory BFGS (L-BFGS) method addresses
the memory limitation with a key observation: we never need to \emph{store} $H_k$
explicitly. What we actually need at each iteration is the matrix-vector product
$H_k \nabla f_k$. If we expand the BFGS recurrence~\eqref{eq:bfgs} backwards
over the last $\bar{m}$ steps, starting from an initial approximation
$H_k^0 = \gamma_k I$, the product $H_k \nabla f_k$ can be expressed as a
sequence of inner products and vector additions involving only the stored pairs
$\{s_i, y_i\}_{i=k-\bar{m}}^{k-1}$. This leads to the \emph{two-loop recursion}
(Algorithm~\ref{alg:twoloop}), which computes $H_k \nabla f_k$ implicitly in
$O(\bar{m}\,m)$ operations, without ever forming or storing $H_k$.

\noindent In practice, $\bar{m} \in [3, 20]$ pairs are retained: Nocedal and
Wright~\cite[Section~9.1]{nocedal2006numerical} note that modest values of
$\bar{m}$ often give satisfactory results, since curvature information from
earlier iterations becomes increasingly irrelevant as the iterates move away from
those regions. At each iteration, the oldest pair is discarded and the newest
pair $(s_k, y_k)$ is added.


% ─────────────────────────────────────────────────────────────
\subsubsection{Initial Scaling of \texorpdfstring{$H_k^0$}{H0}}
\label{sec:h0_scaling}
% ─────────────────────────────────────────────────────────────

\noindent At each iteration, the initial Hessian approximation is set to
$H_k^0 = \gamma_k I$ for some scalar $\gamma_k > 0$. Conceptually, this is the
``blank slate'' from which the two-loop recursion builds $H_k$ by applying the
last $\bar{m}$ BFGS updates. The choice of $\gamma_k$ matters: if it is too
small, the search direction $p_k$ will be short and many small steps will be
needed; if it is too large, the algorithm will overshoot. We implement three
strategies.

% \paragraph{Nocedal scaling (BB2).}
% The default choice~\cite[Eq.~9.6]{nocedal2006numerical} is
% \begin{equation}
% \label{eq:gamma_nocedal}
%     \gamma_k^{\mathrm{N}} = \frac{s_{k-1}^T y_{k-1}}{y_{k-1}^T y_{k-1}}.
% \end{equation}
% This is the scalar $\gamma$ that minimises $\|\gamma y_{k-1} - s_{k-1}\|$,
% i.e.\ the best scalar approximation to the inverse Hessian along the most recent
% step direction $s_{k-1}$. It ensures that the search direction $p_k$ is
% well-scaled and that the unit step length $\alpha_k = 1$ is accepted in most
% Wolfe line-search iterations. This choice coincides with the BB2 step of
% Barzilai and Borwein~\cite{barzilai1988}.


\paragraph{Nocedal scaling (BB2).}
The default choice~\cite[Eq.~9.6]{nocedal2006numerical} is
\begin{equation}
\label{eq:gamma_nocedal}
    \gamma_k^{\mathrm{N}} = \frac{s_{k-1}^T y_{k-1}}{y_{k-1}^T y_{k-1}}.
\end{equation}
This is the unique minimiser of $\|\gamma\, y_{k-1} - s_{k-1}\|^2$ over
$\gamma > 0$, i.e.\ the best scalar approximation to the inverse Hessian along
$s_{k-1}$ in a least-squares sense: $\gamma^{\mathrm{N}}_k$ is the scalar $\gamma$
such that $\gamma y_{k-1} \approx s_{k-1}$, mimicking the secant condition
$H_k^{-1} y_{k-1} = s_{k-1}$.
It ensures that the search direction $p_k$ is well-scaled and that the unit step
length $\alpha_k = 1$ is accepted in most Wolfe line-search iterations.
We note that this formula is algebraically identical to the \emph{BB2} step size
of Barzilai and Borwein~\cite[Eq.~(5)]{barzilai1988}, which was derived
independently in the context of gradient descent by solving the same
least-squares approximation to the secant equation, with
$\Delta x \leftrightarrow s_{k-1}$ and $\Delta g \leftrightarrow y_{k-1}$.
This observation motivates us to also consider the companion BB1 step as an
alternative scaling strategy for $H_k^0$.

\paragraph{BB1 scaling.}
The companion step size proposed by Barzilai and Borwein~\cite[Eq.~(6)]{barzilai1988} is
\begin{equation}
\label{eq:gamma_bb1}
    \gamma_k^{\mathrm{BB1}} = \frac{s_{k-1}^T s_{k-1}}{s_{k-1}^T y_{k-1}}.
\end{equation}
This is the minimiser of $\|\gamma^{-1} s_{k-1} - y_{k-1}\|^2$, i.e.\ the
scalar $\gamma$ such that $\gamma^{-1} s_{k-1} \approx y_{k-1} \approx H\, s_{k-1}$,
yielding an estimate of the \emph{average} eigenvalue of $\nabla^2 f$ (rather
than its inverse). On ill-conditioned problems, BB1 tends to produce larger steps
than BB2, which can accelerate convergence by escaping flat regions more quickly.

% \paragraph{Safeguarded BB2.}
% Following Dai and Liao~\cite{dai2002}, we also provide a clipped variant:
% \begin{equation}
% \label{eq:gamma_safe}
%     \gamma_k^{\mathrm{safe}} = \mathrm{clip}\!\left(
%         \gamma_k^{\mathrm{N}},\; \gamma_{\min},\; \gamma_{\max}
%     \right),
% \end{equation}
% with defaults $\gamma_{\min} = 10^{-10}$ and $\gamma_{\max} = 10^{10}$.
% Near the solution, $y_{k-1}^T y_{k-1}$ can become extremely small (the gradient
% barely changes between steps), making $\gamma_k^{\mathrm{N}}$ blow up and
% rendering the direction $p_k$ unreliable. The clipping prevents this numerical
% instability while retaining the BB2 estimate in the well-behaved regime.


\paragraph{Safeguarded BB2.}
The Nocedal scaling~\eqref{eq:gamma_nocedal} can become numerically unreliable near the 
solution: as $w_k \to w^*$, the gradient changes between successive 
iterates vanish, so $y_{k-1}^T y_{k-1} \to 0$ while $s_{k-1}^T y_{k-1}$ 
remains bounded away from zero, causing $\gamma_k^{\mathrm{N}} \to \infty$.
An unbounded scaling factor renders the search direction $p_k$ unreliable 
regardless of the quality of the two-loop recursion. To prevent this, we 
clip the scaling factor to a finite interval:
\begin{equation}
\label{eq:gamma_safe}
    \gamma_k^{\mathrm{safe}} = \mathrm{clip}\!\left(
        \gamma_k^{\mathrm{N}},\; \gamma_{\min},\; \gamma_{\max}
    \right),
\end{equation}
with defaults $\gamma_{\min} = 10^{-10}$ and $\gamma_{\max} = 10^{10}$.
This retains the BB2 estimate in the well-behaved regime while preventing 
numerical overflow in degenerate cases.

% ─────────────────────────────────────────────────────────────
\subsubsection{Two-Loop Recursion}
% ─────────────────────────────────────────────────────────────

\noindent The product $r = H_k \nabla f_k$ is computed by
Algorithm~\ref{alg:twoloop}~\cite[Algorithm~9.1]{nocedal2006numerical}. The
algorithm implements the BFGS recurrence~\eqref{eq:bfgs} unrolled over the last
$\bar{m}$ pairs, applied implicitly to the gradient vector $q$. The
\emph{backward loop} computes the $\alpha_i$ scalars that will be needed later,
while simultaneously applying the $V_i^T$ factors to $q$ from right to left. The
\emph{scaling step} applies the base approximation $H_k^0 = \gamma_k I$. The
\emph{forward loop} then applies the rank-one corrections $\rho_i s_i s_i^T$,
using the $\alpha_i$ computed in the backward pass. Together, the two loops
implement exactly the product $H_k q$ without ever constructing the matrix.

\begin{algorithm}[H]
\caption{L-BFGS Two-Loop
  Recursion~\cite[Algorithm~9.1]{nocedal2006numerical}}
\label{alg:twoloop}
\begin{algorithmic}[1]
    \Require Gradient $q \gets \nabla f_k$;\;
             stored pairs $\{s_i, y_i\}_{i=k-\bar{m}}^{k-1}$;\;
             scaling $\gamma_k$
    \Ensure  $r = H_k \nabla f_k$
    \For{$i = k-1,\, k-2,\, \ldots,\, k-\bar{m}$}
        \State $\rho_i \gets 1/(y_i^T s_i)$;\quad
               $\alpha_i \gets \rho_i\, s_i^T q$;\quad
               $q \gets q - \alpha_i\, y_i$
        \Comment{Apply $V_i^T$ to $q$}
    \EndFor
    \State $r \gets \gamma_k\, q$
           \Comment{Apply $H_k^0 = \gamma_k I$}
    \For{$i = k-\bar{m},\, \ldots,\, k-1$}
        \State $\beta \gets \rho_i\, y_i^T r$;\quad
               $r \gets r + s_i\,(\alpha_i - \beta)$
        \Comment{Apply rank-one correction}
    \EndFor
    \State \Return $r$
\end{algorithmic}
\end{algorithm}

\noindent The cost of Algorithm~\ref{alg:twoloop} is $4\bar{m}$ dot products of
length $m$ (two in each loop iteration) plus $m$ multiplications for the scaling
step, totalling $(4\bar{m}+1)\,m$ floating-point operations. This is a dramatic
improvement over the $O(m^2)$ cost of storing and applying the full BFGS matrix.


% ─────────────────────────────────────────────────────────────
\subsubsection{Curvature Safeguard and Automatic Restart}
\label{sec:restart}
% ─────────────────────────────────────────────────────────────

\noindent For the BFGS formula~\eqref{eq:bfgs} to produce a positive definite
update, the \emph{curvature condition} $y_k^T s_k > 0$ must hold. To see why,
note that $\rho_k = 1/(y_k^T s_k)$ appears in the rank-one term
$\rho_k s_k s_k^T$: if $y_k^T s_k \leq 0$, then $\rho_k \leq 0$ and the update
would subtract from $H_k$, potentially making $H_{k+1}$ indefinite and the
next direction $p_{k+1}$ a direction of ascent. We implement two levels of
protection against this.

% \paragraph{Negative-curvature reset.}
% If $y_k^T s_k \leq 10^{-16}$, the entire stored memory is cleared and the
% algorithm restarts from a pure gradient-descent step (effectively setting
% $H_k = I$). This is the standard safeguard described in~\cite{liu1989limited}.

% \paragraph{Curvature-based restart.}
% A more anticipatory criterion is proposed in~\cite[Section~4]{liu1989limited}:
% even when $y_k^T s_k > 0$, the memory should be cleared if the curvature
% estimate captured by the new pair changes dramatically relative to the previous
% one. To detect this, we track the quantity
% \begin{equation}
% \label{eq:gamma_restart}
%     \gamma_k = \frac{y_k^T s_k}{y_k^T y_k},
% \end{equation}
% which is the same as the Nocedal scaling~\eqref{eq:gamma_nocedal}: it estimates
% the effective inverse-Hessian scale along $s_k$. A sudden drop,
% \begin{equation}
% \label{eq:restart_cond}
%     \gamma_k < \xi\, \gamma_{k-1}, \qquad \xi \in (0,1),
% \end{equation}
% signals that the curvature along the new step is much smaller than along the
% previous one. This typically happens when the iterates have moved into a
% qualitatively different region of the landscape: the pairs stored in memory
% were computed under a different curvature regime and now give a misleading
% approximation of $H_k^{-1}$. Clearing the memory and restarting from
% $H_k^0 = \gamma_k I$ is then preferable to letting corrupted pairs contaminate
% the two-loop recursion. The threshold $\xi = 0.2$ is used by default.

\paragraph{Negative-curvature reset.}
If $y_k^T s_k \leq 10^{-16}$, the entire stored memory is cleared and the
algorithm restarts from a pure gradient-descent step (effectively setting
$H_k = I$). This safeguard ensures that $\rho_k = 1/(y_k^T s_k)$ remains
well-defined and positive, as required by the BFGS update~\eqref{eq:bfgs}.
The threshold $10^{-16}$ corresponds to the unit roundoff in
double precision arithmetic, below which $\rho_k$ is no longer numerically
reliable.

\paragraph{Curvature-based restart.}
As an additional safeguard, we clear the memory even when $y_k^T s_k > 0$
if the curvature estimate captured by the new pair changes dramatically
relative to the previous one. To detect this, we track the quantity
\begin{equation}
\label{eq:gamma_restart}
    \gamma_k = \frac{y_k^T s_k}{y_k^T y_k},
\end{equation}
which is the same as the Nocedal scaling~\eqref{eq:gamma_nocedal}: it estimates
the effective inverse-Hessian scale along $s_k$. A sudden drop,
\begin{equation}
\label{eq:restart_cond}
    \gamma_k < \xi\, \gamma_{k-1}, \qquad \xi \in (0,1),
\end{equation}
signals that the curvature along the new step is much smaller than along the
previous one. This typically happens when the iterates have moved into a
qualitatively different region of the landscape: the pairs stored in memory
were computed under a different curvature regime and now give a misleading
approximation of $H_k^{-1}$. Clearing the memory and restarting from
$H_k^0 = \gamma_k I$ is then preferable to letting corrupted pairs contaminate
the two-loop recursion. The threshold $\xi = 0.2$ is used by default.


% ─────────────────────────────────────────────────────────────
\subsection{Line Search Strategies}
\label{sec:linesearch}
% ─────────────────────────────────────────────────────────────

\subsubsection{Strong Wolfe Conditions}

\noindent Once the descent direction $p_k$ is computed, the step length $\alpha_k$
must be chosen to make sufficient progress without overshoot. The standard
criterion for quasi-Newton methods is the \emph{strong Wolfe
conditions}~\cite[Eq.~3.7]{nocedal2006numerical}:
\begin{align}
    f(w_k + \alpha_k p_k)
        &\leq f(w_k) + c_1\,\alpha_k\,\nabla f_k^T p_k,
    \label{eq:armijo} \\
    \left|\nabla f(w_k + \alpha_k p_k)^T p_k\right|
        &\leq c_2\,\left|\nabla f_k^T p_k\right|.
    \label{eq:curvature}
\end{align}
Condition~\eqref{eq:armijo} (Armijo) ensures sufficient decrease: the function
must drop by at least a fraction $c_1$ of the directional derivative. Without
it, the algorithm might take steps so small that progress is negligible.
Condition~\eqref{eq:curvature} prevents the step from being too short: the
directional derivative at the new point must be small enough in absolute value,
which ensures that the new pair $(s_k, y_k)$ carries meaningful curvature
information (in particular, it guarantees $y_k^T s_k > 0$).
Following standard practice for quasi-Newton
methods~\cite{nocedal2006numerical}, we set $c_1 = 10^{-4}$ and $c_2 = 0.9$.
Our implementation uses a bracketing-and-zoom procedure with cubic
interpolation~\cite[Section~3.5]{nocedal2006numerical}.


\subsubsection{Exact Line Search for Quadratic Objectives}
\label{sec:exact_ls}

\noindent The Wolfe conditions are designed for general nonlinear objectives and
accept a range of step lengths rather than requiring the exact minimiser. For our
quadratic $f$, we can do much better. The restriction of $f$ along any direction
$p$ is a one-dimensional parabola:
\begin{equation}
\label{eq:phi_parabola}
    \varphi(\alpha) = f(w + \alpha\,p)
    = f(w) + \alpha\,\nabla f(w)^T p + \frac{\alpha^2}{2}\,p^T H p.
\end{equation}
Since $H \succ 0$, the coefficient $p^T H p > 0$ for all $p \neq 0$, so
$\varphi$ is a strictly convex parabola with a unique global minimum at
$\varphi'(\alpha^*) = 0$, i.e.:
\begin{equation}
\label{eq:exact_alpha}
    \alpha^* = -\frac{\nabla f(w)^T p}{\,p^T H p\,}.
\end{equation}
For our specific Hessian $H = XX^T + \lambda^2 I_m$, the denominator can be
expanded without forming $H$ explicitly:
\begin{equation}
    p^T H p = p^T (XX^T + \lambda^2 I_m)\,p
    = \norm{X^T p}^2 + \lambda^2\norm{p}^2,
\end{equation}
so the complete formula becomes a generalisation
of~\cite[Eq.~3.25]{nocedal2006numerical} to an arbitrary descent direction:
\begin{equation}
\label{eq:exact_alpha_expanded}
    \alpha^* = -\frac{\nabla f(w)^T p}{\norm{X^T p}^2 + \lambda^2\norm{p}^2}.
\end{equation}
The only non-trivial cost is the single matrix-vector product $X^T p \in O(mn)$,
identical to the cost of one gradient evaluation, and far cheaper than the
multiple function and gradient evaluations required by the Wolfe bracketing
procedure. Since $\alpha^*$ is the \emph{exact} minimiser of $\varphi$, it fully
exploits the parabolic structure and yields the best possible step at each
iteration, reducing the number of iterations required to reach a given
accuracy compared to the Wolfe variant.


% ─────────────────────────────────────────────────────────────
\subsection{Complete Algorithm}
% ─────────────────────────────────────────────────────────────

\begin{algorithm}[H]
\caption{Enhanced L-BFGS~\cite[Algorithm~9.2]{nocedal2006numerical}}
\label{alg:lbfgs}
\begin{algorithmic}[1]
    \Require Starting point $w_0$; memory $\bar{m}$; tolerance $\varepsilon$;
             scaling strategy $\in \{\mathrm{BB2},\,\mathrm{BB1},\,\mathrm{safe}\}$;
             restart flag; restart threshold $\xi$
    \State $k \gets 0$;\; compute $\nabla f_0$;\;
           $\tau \gets \varepsilon\,\|\nabla f_0\|$;\;
           $\gamma_{\mathrm{prev}} \gets 1$
    \While{$\|\nabla f_k\| > \tau$}
        \State Compute $\gamma_k$ via chosen scaling strategy
               (Eq.~\eqref{eq:gamma_nocedal},~\eqref{eq:gamma_bb1},
               or~\eqref{eq:gamma_safe})
        \State $p_k \gets -H_k \nabla f_k$
               \Comment{Via Algorithm~\ref{alg:twoloop} with $H_k^0 = \gamma_k I$}
        \If{$\nabla f_k^T p_k \geq 0$}
            \State $p_k \gets -\nabla f_k$
            \Comment{Descent safeguard: fall back to steepest descent}
        \EndIf
        \State Compute $\alpha_k$ via exact LS~\eqref{eq:exact_alpha_expanded}
               or Wolfe conditions~\eqref{eq:armijo}--\eqref{eq:curvature}
        \State $w_{k+1} \gets w_k + \alpha_k p_k$;\quad
               $s_k \gets \alpha_k p_k$;\quad
               $y_k \gets \nabla f_{k+1} - \nabla f_k$

        \State $\gamma_k \gets (y_k^T s_k)/(y_k^T y_k)$
        \If{$y_k^T s_k > 10^{-10}$}
            \If{restart enabled \textbf{and}
                $\gamma_k < \xi\,\gamma_{\mathrm{prev}}$}
                \State Clear all stored pairs
                \Comment{Curvature-based restart, Eq.~\eqref{eq:restart_cond}}
            \EndIf
            \State Store $(s_k, y_k)$;\; discard oldest if
                   $\#\text{stored} > \bar{m}$;\;
                   $\gamma_{\mathrm{prev}} \gets \gamma_k$
                   
        \algstore{lbfgs_break}       
\end{algorithmic}
\end{algorithm}

\begin{algorithm}[H]
\label{alg:lbfgs_part2}
\begin{algorithmic}[1]
        \algrestore{lbfgs_break}
    
        \ElsIf{$y_k^T s_k \leq 10^{-16}$}
            \State Clear all stored pairs
            \Comment{Negative-curvature reset}
        \Else
            \State Discard $(s_k, y_k)$
            \Comment{Curvature too small to be reliable; pair skipped}
        \EndIf
        \State $k \gets k + 1$
    \EndWhile
    \State \Return $w_k$
\end{algorithmic}
\end{algorithm}


% ─────────────────────────────────────────────────────────────
\subsection{Stopping Criterion}
\label{sec:stopping}
% ─────────────────────────────────────────────────────────────

\noindent The primary stopping criterion uses a \emph{relative} gradient tolerance:
\begin{equation}
\label{eq:stop_rel}
    \|\nabla f(w_k)\| \leq \varepsilon \cdot \|\nabla f(w_0)\|.
\end{equation}
The relative form is preferred over an absolute threshold $\|\nabla f(w_k)\| \leq
\varepsilon$ because it is scale-invariant: the magnitude of $\nabla f$ depends
on the problem scale (and in particular on $\lambda$), so an absolute threshold
that is appropriate for one value of $\lambda$ may be either too loose or too
tight for another. Dividing by $\|\nabla f(w_0)\|$ normalizes for this effect.

\noindent As a secondary safeguard, the algorithm detects \emph{numerical
stagnation}: if the relative change $|f_k - f_{k-1}|/|f_k|$ drops below machine
epsilon (approximately $10^{-16}$) while the step length $\alpha_k \approx 0$,
further progress is impossible due to floating-point limitations, and the
iteration is terminated early.


% % ─────────────────────────────────────────────────────────────
% \subsection{Computational Complexity}
% \label{sec:complexity}
% % ─────────────────────────────────────────────────────────────

% \noindent Table~\ref{tab:flops} breaks down the dominant floating-point operations
% per iteration, derived directly from the structure of
% Algorithm~\ref{alg:twoloop} and the quadratic objective. Throughout, we assume
% the exact line search is used, which requires a single $X^T p$ product.

% \begin{table}[h]
% \centering
% \caption{Flop budget per L-BFGS iteration (exact line search).}
% \label{tab:flops}
% \begin{tabular}{llr}
%     \toprule
%     \textbf{Operation}          & \textbf{Asymptotic cost} & \textbf{Exact count} \\
%     \midrule
%     Gradient evaluation         & $O(mn)$              & $2mn$ \\
%     Two-loop recursion          & $O(\bar{m}\,m)$      & $(4\bar{m}+1)\,m$ \\
%     Exact line search ($X^Tp$)  & $O(mn)$              & $mn + 2m$ \\
%     Pair update ($s_k,\,y_k$)   & $O(m)$               & $2m$ \\
%     \midrule
%     \textbf{Total per iteration}& $O(mn + \bar{m}\,m)$ & $m(4\bar{m} + 3n + 5)$ \\
%     \bottomrule
% \end{tabular}
% \end{table}

% \noindent The memory requirement for the stored pairs is $(2\bar{m}+4)\,m$
% floating-point numbers: $2\bar{m}$ vectors of length $m$ for the pairs
% $\{s_i, y_i\}$, plus four more for the current $w$, $\nabla f$, $s_k$, $y_k$.
% Table~\ref{tab:cost_mem} quantifies cost and storage for the ML-CUP problem.

% \begin{table}[h]
% \centering
% \caption{Theoretical cost and storage as a function of memory size
%   ($m=500$, $n=12$, $K=50$ iterations).}
% \label{tab:cost_mem}
% \begin{tabular}{rrrr}
%     \toprule
%     $\bar{m}$ & \textbf{Flops/iter} & \textbf{Total flops} ($K=50$)
%               & \textbf{Storage (MB)} \\
%     \midrule
%      3  &  26{,}500 &  1{,}325{,}000 & 0.040 \\
%      5  &  30{,}500 &  1{,}525{,}000 & 0.056 \\
%     10  &  40{,}500 &  2{,}025{,}000 & 0.096 \\
%     20  &  60{,}500 &  3{,}025{,}000 & 0.176 \\
%     40  & 100{,}500 &  5{,}025{,}000 & 0.336 \\
%     \bottomrule
% \end{tabular}
% \end{table}

% \noindent Two observations stand out. First, the total cost scales \emph{linearly}
% in $\bar{m}$: doubling the memory doubles the per-iteration cost, but also
% potentially halves the number of iterations required (up to the saturation
% point discussed in Section~\ref{sec:memory_effect}). Second, even at
% $\bar{m} = 40$, the storage for the ML-CUP problem ($m = 500$) is less than
% $0.4$\,MB, negligible compared to the $O(m^2) = 2$\,MB required by full BFGS.
% This contrast becomes even more striking in high-dimensional problems: for
% $m = 10^5$, full BFGS would require ${\sim}80$\,GB, while L-BFGS with
% $\bar{m} = 10$ needs only ${\sim}16$\,MB.


% ─────────────────────────────────────────────────────────────
\subsection{Convergence Analysis}
\label{sec:convergence}
% ─────────────────────────────────────────────────────────────

\subsubsection{Global Linear Convergence}

\noindent The convergence of L-BFGS on our problem follows
from~\cite[Theorem~7.1]{liu1989limited}. That theorem requires three conditions,
which we verify directly using the properties established in
Section~\ref{sec:properties}:
\begin{enumerate}[nosep]
    \item $f \in C^2$: satisfied, since $f$ is quadratic;
    \item the sublevel set $\mathcal{L}_0 = \{w : f(w) \leq f(w_0)\}$ is
          bounded: satisfied, since $f$ is $\mu$-strongly convex with
          $\mu = \lambda^2 > 0$, so all sublevel sets are compact ellipsoids;
    \item there exist $M_1, M_2 > 0$ such that
          $M_1\|v\|^2 \leq v^T \nabla^2 f(w)\,v \leq M_2\|v\|^2$ for all $v$:
          satisfied with $M_1 = \mu = \lambda^2$ and
          $M_2 = L = \lambda_{\max}(XX^T) + \lambda^2$, since
          $\nabla^2 f = H$ is constant.
\end{enumerate}
Under these conditions, the iterates satisfy the linear convergence bound:
\begin{equation}
\label{eq:linconv}
    f(w_k) - f(w^*) \leq r^k\,(f(w_0) - f(w^*)), \quad r \in [0,1),
\end{equation}
where the rate $r$ depends on the condition number $\kappa(\hat{X})$: the more
ill-conditioned the problem, the slower the linear rate.

\subsubsection{Superlinear Convergence}

\noindent For full BFGS (i.e.\ $\bar{m}$ large enough to store the entire
history),~\cite[Theorem~8.6]{nocedal2006numerical} guarantees superlinear
convergence on strongly convex functions: the error ratio
\begin{equation}
\label{eq:superlin}
    \frac{\norm{w_{k+1} - w^*}}{\norm{w_k - w^*}} \longrightarrow 0.
\end{equation}
Intuitively, as $H_k$ accumulates more and more
accurate curvature information, it converges to the true $H^{-1}$, and the
search direction $p_k \approx -H^{-1}\nabla f_k$ becomes increasingly close to
the Newton step.

% \noindent A non-asymptotic analysis by Rodomanov and Nesterov~\cite{rodomanov2021}
% predicts that superlinear behaviour begins after approximately
% $K_0 = 4m\ln\kappa$ iterations. Jin and Mokhtari~\cite{jin2022} provide the
% local bound
% \begin{equation}
%     \frac{\|w_k - w^*\|_H}{\|w_0 - w^*\|_H}
%     \leq \left(\frac{1}{k}\right)^{k/2},
% \end{equation}
% where $\|v\|_H = \sqrt{v^T H v}$ is the energy norm. This implies that
% $\varepsilon$-accuracy requires only
% $k = O\!\left(\frac{\log(1/\varepsilon)}{\log\log(1/\varepsilon)}\right)$
% iterations, substantially better than the $O(\log(1/\varepsilon))$ bound for
% linear convergence.


% % ─────────────────────────────────────────────────────────────
% \subsection{Applicability to the Problem at Hand}
% \label{sec:applicability}
% % ─────────────────────────────────────────────────────────────

% \noindent The convergence results above are stated for general objective functions
% satisfying a set of regularity conditions. We now verify explicitly that all such
% conditions hold for our specific problem, so that the theoretical guarantees
% apply rigorously.

% \begin{enumerate}

%     \item \textbf{Twice differentiability.}
%     The objective $f(w) = \tfrac{1}{2}\|X^T w - y\|^2 + \tfrac{1}{2}\lambda^2\|w\|^2$
%     is a quadratic polynomial in $w$, hence $f \in C^\infty$. In particular,
%     the Hessian $H = XX^T + \lambda^2 I_m$ is constant and globally defined.

%     \item \textbf{Bounded sublevel sets.}
%     Since $f$ is $\mu$-strongly convex with $\mu = \lambda^2 > 0$
%     (Section~\ref{sec:properties}), all sublevel sets
%     $\mathcal{L}_c = \{w : f(w) \leq c\}$ are compact ellipsoids.
%     In particular, the initial sublevel set $\mathcal{L}_0$ is bounded.

%     \item \textbf{Lipschitz continuous gradient.}
%     The gradient $\nabla f(w) = Hw - Xy$ is a linear function of $w$,
%     hence Lipschitz continuous with constant $L = \lambda_{\max}(H)$
%     (Section~\ref{sec:properties}):
%     \begin{equation}
%         \|\nabla f(w) - \nabla f(w')\| = \|H(w - w')\|
%         \leq L\|w - w'\| \quad \forall\, w, w' \in \mathbb{R}^m.
%     \end{equation}

%     \item \textbf{Line search condition.}
%     The exact line search (Section~\ref{sec:exact_ls}) finds the global
%     minimiser of $f$ along each direction and therefore satisfies any
%     reasonable descent condition. The strong Wolfe line search
%     (Section~\ref{sec:linesearch}) satisfies
%     conditions~\eqref{eq:armijo}--\eqref{eq:curvature} by construction, and
%     in particular guarantees $y_k^T s_k > 0$ at every accepted step, ensuring
%     that all BFGS updates are well-defined and positive definite.

% \end{enumerate}

% \noindent All four conditions are satisfied. Therefore, both the global linear
% convergence result of~\cite{liu1989limited} and the local superlinear
% convergence result of~\cite{nocedal2006numerical} apply rigorously to our
% setting, with explicit constants $\mu = \lambda^2$ and
% $L = \lambda_{\max}(XX^T) + \lambda^2$.